\section{平方の差を作る双曲線関数}\label{節:双曲線関数}

円では$X^2+Y^2=1$を使った。これに対して，$\sqrt{1+x^2}$を簡単にしたい場合には，$Y^2-X^2=1$という平方の差を使えればよい。指数関数の和と差が，その形を作る。

\begin{definition}{双曲線関数}
実数$u$に対して
\[
\sinh u=\frac{\e^u-\e^{-u}}2,\qquad
\cosh u=\frac{\e^u+\e^{-u}}2,\qquad
\tanh u=\frac{\sinh u}{\cosh u}
\]
と定義する。順に双曲線正弦，双曲線余弦，双曲線正接という。
\end{definition}

$\cosh u>0$なので$\tanh u$はすべての実数で定義される。和と差の積から
\[
\begin{aligned}
\cosh^2u-\sinh^2u
&=\frac{(\e^u+\e^{-u})^2-(\e^u-\e^{-u})^2}{4}\\
&=\e^u\e^{-u}=1
\end{aligned}
\]
となる。従って$(X,Y)=(\sinh u,\cosh u)$は，双曲線$Y^2-X^2=1$の上側にある。図\ref{図:双曲線の媒介表示}のように，$u$が正なら$X>0$，負なら$X<0$である。この$u$は，円のときのような通常の中心角として定義したものではない。

\begin{figure}[htbp]
\centering
\begin{tikzpicture}[x=1.5cm,y=1.5cm,>=Stealth]
 \draw[->] (-2.4,0)--(2.55,0) node[right]{$X$};
 \draw[->] (0,-0.15)--(0,2.8) node[above]{$Y$};
 \draw[dashed,gray] (-2.2,2.2)--(0,0)--(2.2,2.2);
 \draw[very thick,tealblack,domain=-2.2:2.2,samples=80]
  plot(\x,{sqrt(1+\x*\x)});
 \draw[dashed,blueblack] (1.175201,0)--(1.175201,1.543081)--(0,1.543081);
 \fill[blueblack] (1.175201,1.543081) circle(1.3pt);
 \node[anchor=west] at (0.25,2.55) {$Y^2-X^2=1$};
 \node[anchor=east] at (-0.3,1.65) {$u<0$};
 \node[anchor=west] at (0.35,1.95) {$u>0$};
 \node[anchor=west] at (1.4,0.58) {$(\sinh u,\cosh u)$};
 \draw[->,blueblack] (1.75,0.82)--(1.2,1.48);
 \node[left] at (0,1) {$1$};
 \node[below left] at (0,0) {$0$};
\end{tikzpicture}
\caption{指数関数の和と差によって，双曲線の上側を一度ずつ表す。点は$u=1$の位置。}
\label{図:双曲線の媒介表示}
\end{figure}

定義を微分すると
\[
(\sinh u)'=\frac{\e^u+\e^{-u}}2=\cosh u,
\qquad
(\cosh u)'=\frac{\e^u-\e^{-u}}2=\sinh u.
\]
余弦の場合と違って，$\cosh$を微分しても負号は付かない。商の微分から
\begin{equation}\label{式:双曲線正接の微分}
(\tanh u)'=\frac{\cosh^2u-\sinh^2u}{\cosh^2u}
=\frac1{\cosh^2u}=1-\tanh^2u
\end{equation}
である。加えて指数関数の定義式を展開すると
\begin{equation}\label{式:双曲線の二倍公式}
\sinh(2u)=2\sinh u\cosh u,\qquad
\cosh(2u)=2\cosh^2u-1
\end{equation}
となる。例えば後者の右辺は
\[
\frac{(\e^u+\e^{-u})^2}{2}-1
=\frac{\e^{2u}+\e^{-2u}}2
\]
であり，左辺の定義に一致する。

\subsection{双曲線正弦から元へ戻す}

$\cosh u>0$なので，$\sinh u$は実数全体で狭義単調増加する。また$u\to\infty$で$\sinh u\to\infty$，$u\to-\infty$で$\sinh u\to-\infty$だから，どの実数$x$も$x=\sinh u$と一意に書ける。この逆関数を$\operatorname{arsinh}x$と書く。

$x=\sinh u$なら，$\cosh u=\sqrt{1+x^2}$である。定義式を足すと$\sinh u+\cosh u=\e^u$だから
\begin{equation}\label{式:逆双曲線正弦}
u=\operatorname{arsinh}x=\log\bigl(x+\sqrt{1+x^2}\bigr).
\end{equation}
対数の中は，$\sqrt{1+x^2}>|x|$によって常に正である。逆関数の微分から
\[
\frac{\dd u}{\dd x}=\frac1{\cosh u}
=\frac1{\sqrt{1+x^2}}
\]
となる。ここでも積分の値から逆関数を定義したのではなく，指数関数による定義から順に導いた。

\begin{example}{根号の分母を双曲線正弦で取り除く}
$\int_0^3\frac1{\sqrt{4+x^2}}\dd x$を求めよ。
\end{example}

$x=2\sinh u$と置けば
\[
\dd x=2\cosh u\dd u,\qquad
\sqrt{4+x^2}=2\cosh u.
\]
分母と微分の係数が打ち消し合う。端点$x=0$は$u=0$へ，$x=3$は
\[
u=\operatorname{arsinh}\frac32
=\log\frac{3+\sqrt{13}}2
\]
へ移る。従って積分値は$\log\frac{3+\sqrt{13}}2$である。三角置換$x=2\tan\theta$では$\sec\theta$の積分が残るが，この置換では$1$の積分になる。

$\cosh u$は偶関数なので，実数全体では逆関数を持たない。しかし$u\geqq0$に制限すれば，$1$から無限大へ狭義単調増加する。$x=\cosh u$を使う置換では，このように枝を選ぶ。$\sinh$と$\cosh$で，逆向きに戻すときの条件は同じではない。

\begin{example}{平方の差を双曲線正接で取り除く}
$\int_0^{\frac12}\frac1{1-x^2}\dd x$を求めよ。
\end{example}

$\tanh u=\frac{\e^{2u}-1}{\e^{2u}+1}$は実数全体で狭義単調増加し，値の範囲は$(-1,1)$である。$x=\tanh u$と置けば，式\eqref{式:双曲線正接の微分}によって$\dd x=(1-x^2)\dd u$となる。元へ戻す式は
\[
x(\e^{2u}+1)=\e^{2u}-1,\qquad
\e^{2u}=\frac{1+x}{1-x},\qquad
u=\frac12\log\frac{1+x}{1-x}
\]
である。従って端点は$0$と$\frac12\log3$へ移り，求める積分値は$\frac12\log3$となる。部分分数で
\[
\frac1{1-x^2}=\frac12\left(\frac1{1+x}+\frac1{1-x}\right)
\]
と分けても同じ対数が現れる。双曲線正接の逆関数が，その二つの対数を一つにまとめている。

%
